% Author: Lisa Lusinga Namugenyi (F-03, F-11). Draft written by Charis Opol on 23 Nov 2026 for Lisa
% to check and extend. Every number is from one of:
%   analysis/01_grid/tables/test_wer_greedy.md, test_wer_grammar.md, missing_modality.md
%   analysis/01_grid/bootstrap.txt (paired bootstrap and per-speaker WER;
%     analysis/scripts/grid_bootstrap.py, built from reports/<run>/hypotheses.csv)
%   reports/gate_analysis/summary.md
%   reports/w4_05_curriculum_comparison.md (training budget and epochs)
\section{Results on GRID}
\label{sec:results}

\subsection{Protocol}
Every GRID model is tested once on the four held-out speakers, at each of the five noise levels,
on the fixed noisy copies of the test set (Section~\ref{sec:data}), giving 3,998 scored clips per
level. WER and CER are total edits divided by total reference words or characters over the test
set. Headline numbers use greedy decoding; the grammar decoder is reported beside them. To compare
two models, we resample the test clips with replacement 2,000 times and recompute the difference in
WER on the same clips, which gives a 95~\% paired bootstrap interval
(\texttt{analysis/\allowbreak{}scripts/\allowbreak{}grid\_\allowbreak{}bootstrap.py}). The four test speakers are too few to resample by
speaker, so these intervals cover variation between sentences, not between speakers. The P models
were each trained once (seed~1) within the same 30-minute budget
(\texttt{reports/\allowbreak{}w4\_\allowbreak{}05\_\allowbreak{}curriculum\_\allowbreak{}comparison.md}), and the ablations once each with the same
seed; like P4, each ablation ran for 6 epochs \textbf{[AQ1]}.

\begin{table*}[t]
\caption{GRID test WER (\%) on four unseen speakers, greedy decoding (grammar decoding in
parentheses for the final models). The lips-only models do not hear the noise, so their WER is the
same at every level.}
\label{tab:grid}
\centering
\setlength{\tabcolsep}{5pt}
\begin{tabular}{llrrrrr}
\toprule
Model & Streams and initialisation & clean & 10 dB & 5 dB & 0 dB & $-5$ dB \\
\midrule
B1      & audio only                         & 1.96 (1.76) & 3.02 (2.79) & 3.79 (3.47) & 5.44 (4.93) & 9.45 (8.35) \\
B2 v1   & lips only, 3-D CNN (V1)            & 23.04 & 23.04 & 23.04 & 23.04 & 23.04 \\
B2 v2   & lips only, ResNet-18 (V2)          & 32.20 & 32.20 & 32.20 & 32.20 & 32.20 \\
B3 v1   & concatenation, from scratch        & 2.41 & 2.88 & 3.36 & 4.29 & 7.35 \\
B3 v2   & concatenation, from B1 + B2 v2     & 1.16 (0.77) & 1.66 (1.13) & 1.98 (1.38) & 2.81 (1.99) & 5.21 (3.83) \\
\midrule
P1      & model P, from scratch              & 12.58 & 13.05 & 13.55 & 14.49 & 16.28 \\
P2      & model P, from B2 v1                & 18.19 & 18.20 & 18.53 & 18.93 & 19.85 \\
P3      & model P, from B1 + B2 v1           & 2.13 & 2.67 & 3.27 & 4.33 & 7.30 \\
P4      & model P, from B1 + B2 v2           & \textbf{1.02} (0.71) & \textbf{1.53} (1.10) & \textbf{1.87} (1.42) & \textbf{2.79} (2.06) & 5.34 (3.95) \\
\bottomrule
\end{tabular}
\end{table*}

% Figure made by figures/make_wer_snr.py. The IEEE Access class needs its \Figure macro: a plain
% figure environment fails in \@makecaption (\xfigwd undefined). Full width because the image is
% wider than one column. The class sets a wide figure's caption in a one-line l column, so the
% caption text is wrapped in a \parbox to let it break.
\Figure[t!](topskip=0pt, botskip=0pt, midskip=0pt)[width=\textwidth]{figures/wer_vs_snr.pdf}
{\parbox[t]{0.9\textwidth}{Test WER against noise level for the audio-only model (dashed) and the
audio + lips model (solid) on each source. GRID: B1 and P4, greedy. TED and films: twin\_v5 and
ws\_v5, beam search with the language model (Table~\ref{tab:realworld}). Absolute WERs are not
comparable across panels, because GRID has a 51-word vocabulary and films and TED are open
vocabulary; the gap within each panel is the quantity of interest.}\label{fig:wersnr}}

\subsection{Fusion against the audio-only baseline}
Table~\ref{tab:grid} and Fig.~\ref{fig:wersnr} answer the first question. P4 makes about half the
errors of the audio-only B1 at every level: 1.02~\% against 1.96~\% on clean audio and 5.34~\%
against 9.45~\% at $-5$~dB. The relative reduction stays between 43.5 and 50.7~\%, so the absolute
gain grows as the noise rises, from 0.94 points on clean audio (95~\% interval 0.77 to 1.11) to
4.11 points at $-5$~dB (3.74 to 4.49). Every interval excludes zero. In terms of effective SNR
gain~\cite{lin2025}, P4 at $-5$~dB (5.34~\%) is still better than B1 at 0~dB (5.44~\%), so the
lips are worth more than 5~dB of noise on GRID; we cannot say how much more, because no level below
$-5$~dB was tested. The grammar decoder lowers P4's WER further, to 0.71~\% clean and 3.95~\% at
$-5$~dB, well inside the proposal's target of less than 10~\% on clean audio.

The gain holds for each test speaker in noise. At $-5$~dB, P4 lowers the WER of all four speakers,
from 7.60 to 5.17~\% (s14), 6.85 to 3.35~\% (s24), 9.40 to 3.76~\% (s4) and 13.97 to 9.10~\%
(s9). On clean audio the picture is mixed: s4 and s9 gain 2.00 and 1.87 points, s24 ties at
0.30~\%, and s14 is slightly worse (1.48~\% against 1.38~\%).

\subsection{Effect of initialisation and fusion type}
Within the 30-minute budget, model P learned well only when both encoders started from trained
baselines (Table~\ref{tab:grid}). From scratch (P1), or with only the lip encoder pretrained (P2),
it stayed above 12~\% WER even on clean audio. Starting from B1 and a lip encoder (P3, P4) brought it
below the audio-only baseline. P4, with the ResNet-18 encoder, beat P3 at every level, even though
its lips-only starting point B2~v2 was weaker on its own than B2~v1 (32.20~\% against 23.04~\%).

The concatenation baseline B3~v2, which starts from the same two encoders and also uses modality
dropout, comes close to P4. P4 is 0.14 points better on clean audio (interval 0.01 to 0.27), but
at 10, 5 and 0~dB the intervals include zero, and at $-5$~dB B3~v2 is 0.14 points better, with an
interval ($-0.12$ to $0.41$) that also includes zero. Most of P4's advantage over B1 therefore comes
from adding a well-initialised lip stream with modality dropout, not from attention and gating on
their own. Ablation X2 below sharpens this point.

\subsection{Ablations}
Table~\ref{tab:ablation} removes one part of P4 at a time (Section~\ref{sec:ablations}). Every
change makes P4 worse at every noise level, and every interval excludes zero.

\begin{table}[t]
\caption{Ablations of P4 on the GRID test set: WER (\%), greedy, and the change from P4 with its
95~\% paired bootstrap interval. Positive means the ablation is worse.}
\label{tab:ablation}
\centering
\setlength{\tabcolsep}{1.8pt}
\begin{tabular}{lrrrr}
\toprule
SNR & P4 & X1 (no dropout) & X2 (concat) & X3 (2-D) \\
\midrule
clean   & 1.02 & 1.57 & 1.23 & 1.40 \\
        &      & $+0.55$ [0.42, 0.70] & $+0.20$ [0.08, 0.33] & $+0.38$ [0.25, 0.52] \\
10 dB   & 1.53 & 2.09 & 1.78 & 2.22 \\
        &      & $+0.55$ [0.40, 0.71] & $+0.25$ [0.10, 0.40] & $+0.69$ [0.53, 0.85] \\
5 dB    & 1.87 & 2.56 & 2.32 & 2.63 \\
        &      & $+0.70$ [0.53, 0.88] & $+0.45$ [0.29, 0.62] & $+0.77$ [0.58, 0.95] \\
0 dB    & 2.79 & 3.54 & 3.19 & 3.92 \\
        &      & $+0.74$ [0.53, 0.95] & $+0.40$ [0.19, 0.60] & $+1.13$ [0.90, 1.36] \\
$-5$ dB & 5.34 & 6.33 & 5.89 & 7.50 \\
        &      & $+0.98$ [0.70, 1.28] & $+0.55$ [0.26, 0.83] & $+2.16$ [1.85, 2.48] \\
\bottomrule
\end{tabular}
\end{table}

\begin{itemize}
\item \textit{Modality dropout (X1)} is worth 0.55 points on clean audio and 0.98 points at
$-5$~dB. A model that always sees both streams learns to work less well from either.
\item \textit{Attention and gating (X2)} are worth 0.20 to 0.55 points. They help, but they are the
smallest of the three parts, which agrees with the closeness of B3~v2 and P4. X2 is also worse
than B3~v2, although both concatenate the streams; the two differ in their projection layers and
training length (7 epochs for B3~v2, 6 for X2), and we did not isolate which difference matters
\textbf{[AQ2]}.
\item \textit{The 3-D front end (X3)} matters most in noise: its loss grows from 0.38 points on
clean audio to 2.16 points at $-5$~dB, a 40~\% relative increase in WER. When the model leans on
the lips, the short-range lip motion captured by the temporal kernel carries much of the useful
information.
\end{itemize}

\subsection{Missing modality}
With the video removed at test time and replaced by $e_{\text{miss}}$, P4 behaves like an audio-only
model: 1.88~\% on clean audio (B1: 1.96~\%) and 10.36~\% at $-5$~dB (B1: 9.45~\%)
(\texttt{analysis/\allowbreak{}01\_\allowbreak{}grid/\allowbreak{}tables/\allowbreak{}missing\_\allowbreak{}modality.md}). With the audio removed, it reads the lips
at 31.57~\%, close to its lip encoder's starting point B2~v2 (32.20~\%). The concatenation model
B3~v2, whose missing stream is replaced by zeros, falls apart instead: 75.31~\% from the lips alone
and 25.18~\% from clean audio alone. The learned $e_{\text{miss}}$ embedding thus allows P4 to degrade gracefully when one stream is
unavailable, as can happen in live use.

\subsection{Gate behaviour}
Fig.~\ref{fig:gate} shows P4's gate, averaged over the frames of each test clip. Its mean falls
steadily from 0.804 on clean audio to 0.786, 0.769, 0.740 and 0.691 at 10, 5, 0 and $-5$~dB, and its
spread widens, so in noise some clips lean much more on the lips than others. The Spearman
correlation between SNR and the per-clip gate is 0.47 (\texttt{reports/\allowbreak{}gate\_\allowbreak{}analysis/\allowbreak{}summary.md}).
The gate never favours the lips on average: even at $-5$~dB about two thirds of the weight goes to
the audio. The gate therefore shifts towards the lips as noise increases, but the audio remains the
dominant stream at every tested level.

% Width just under \columnwidth, so \Figure keeps it in one column.
\Figure[t!](topskip=0pt, botskip=0pt, midskip=0pt)[width=0.98\columnwidth]{figures/gate_boxplot.png}
{P4's gate per test clip (mean over frames) at each noise level; 1 means full trust in the
audio, 0 full trust in the lips. Triangles mark the mean. 3,998 clips per level.\label{fig:gate}}

\subsection{Comparison with published GRID results}
The closest published numbers are not directly comparable, because the speaker splits and the
tasks differ. Our best lips-only model, B2~v1, has a WER of 23.04~\% on four unseen speakers, which
is of the same order as the 79.6~\% word accuracy of Wand et al.\ on their held-out
speakers~\cite{wand2016}. LipNet's 95.2~\% sentence-level accuracy and LCANet's 3.0~\%
WER~\cite{assael2016,xu2018} come from splits in which every test speaker also appears in training
(LCANet holds out 255 sentences of each speaker), which is an easier task than ours.
